\section{Introduction}
A renewal service provider may observe many histories while its terminal gateway accepts only a small menu of commands. The provider must choose that menu and allocate an accepted service obligation across histories. Participation limits expected service on each history, a separate ceiling limits every realized delivery, and installing a command can incur a fixed charge. These constraints make the menu and the promised service targets interdependent. Optimizing commands for prescribed targets does not solve the joint design problem.

We develop a joint-design method that preserves the original command budget, charges, aggregate obligation, and individual constraints. Its first component is an exact price oracle on a box of target intervals. A mandatory lower target can make commands after the price-optimal peak indispensable. An increasing prefix accounts for upper-target crossings, a decreasing suffix accounts for lower-target crossings, and their common peak joins two path recurrences. The oracle gives valid upper bounds over every menu in the box, not just over a menu supplied by a heuristic. Adaptive partitioning then yields an interval containing the global joint optimum, with finite termination at every positive additive tolerance.

The second component reduces the history dimension without requiring identical histories. Groups retain the same eligible commands but may differ in caps and costs. Weighted caps preserve aggregate capacity, while a minimum-cap guard preserves feasible commands. We strengthen minimum-curvature aggregation by using weighted harmonic service curvatures. This gives a provably tighter joint upper model without extra target coordinates. Recovery then optimizes each group's original allocation at its reduced mean, preserving the installed menu, exact obligation, individual caps, and every realization ceiling.

A square-root curvature partition gives a second-order cost-dispersion bound, including zero costs. Together with cap-dispersion control, it improves the accuracy-dependent dimension of the variable-history additive scheme. For bounded eligibility complexity and primitive scales, fixed-accuracy complexity is polynomial in the number of histories and catalog size. The accuracy dependence can still be large. Practical coarse partitions are judged by a checked original-space interval, never by the reduced solver's completion flag alone. Refinement and the best available original-space witnesses provide a constructive route to tighter certificates.

An independent checker verifies the group construction, recovery identities, original lotteries, price-path inequalities, and the complete target cover. The computational evidence includes distinct histories and a paired cost-heterogeneity ablation of harmonic and minimum-curvature aggregation. It compares certified coarsening with unaggregated adaptive search, the inherited priced-prefix method, exhaustive menu enumeration, and an uneliminated mixed-integer formulation. All declared cases and unresolved intervals are retained, including regimes where cap dispersion dominates the improved cost bound. Separate stress instances show that an exact reduced solution can still leave a substantial original-space gap, and that refinement removes this gap. These synthetic experiments evaluate algorithms and certificates, not an estimated operational population.

\paragraph{Relation to established methods.}\label{sec:r37-position}
Classical interval quantization and matrix searching underpin the saturated and conditional path algorithms \citep{Wu1991,NielsenNock2014,GronlundEtAl2017}. Pricing, constrained paths, and branch-and-bound are established tools \citep{HandlerZang1980,LawlerWood1966,Lemarechal2001}. Spatial branch-and-bound also exploits locally tight relaxations of separable piecewise-linear problems \citep{HubnerEtAl2025}. We do not claim those general mechanisms as new. Our two-sided oracle handles a shared, charged, cardinality-constrained menu under heterogeneous eligibility, with an anchor condition that ensures local tightness.

Aggregation error analysis is likewise established. \citet{LiBertsekas2025} bound value-function error in aggregated discounted dynamic programming, and \citet{Fattahi2025} studies approximation and error components in jointly designing charging programs and managing a heterogeneous participant population. Our result concerns a different coupling: one installed menu and an exact root obligation must survive aggregation and recovery while every original cap and realization ceiling remains enforceable. The operative ingredients are eligibility-preserving groups, the minimum-cap guard, and a group-mass-preserving lift. The inequalities between arithmetic and harmonic means are classical; see, for example, \citet{LiaoWu2015}. Their use here must additionally preserve the shared menu and every original contract constraint. The associated complexity guarantee does not assume a small number of exactly repeated types.

The broader contractual and continuous-design contributions remain part of the paper. Canonical allocation connects expected and pathwise participation; continuous quadratic design, Monge acceleration, conditional compression, full target nets, and charge-aware recovery retain their original assumptions and proofs. Table~\ref{tab:r46-contributions} separates these antecedents from two-sided joint decomposition and certified response coarsening. The article presents the joint-design argument first; the companion retains the full preceding mathematical chain. The repository computational record carries the intact earlier study narratives and the complete new per-instance measurements.
